\clearpage
\setcounter{page}{10}
\begin{absolutelynopagebreak}
\chapter{Abstract}

The Kitaev honeycomb model realizes a quantum spin liquid where spin degrees of freedom are fractionalized into itinerant Majorana fermions and localized vortices.
The entanglement among these particles holds promise for executing fault-tolerant topological quantum computation.
While several materials have been proposed to host a quantum spin liquid phase, concrete evidence confirming its existence remains elusive.
To aid in identifying this phase, we propose an electrical test using the charge profile, offering a signature for the $\mathbb{Z}_2$ vortices in Kitaev materials.


This study explores the electric charge response in Kitaev materials, considering a generalized spin model including Heisenberg, off-diagonal exchange interactions, and magnetic Zeeman field.
Majorana mean-field theory is employed to analyze spin correlations around vortices, revealing spatial anisotropy in the charge distribution, quantified by the electric quadrupole moment.
This mean-field approach remains exact in the pure Kitaev limit and captures the fundamental physics across the Kitaev spin liquid phase.
The quadrupole-quadrupole interaction between distant vortices is shown to be either repulsive or attractive, depending on the parameters. We predict that electrically biased scanning probe tips can help create vortices at predetermined locations.

The findings of our study pave the way for electrically manipulating Ising anyons in Kitaev spin liquids, offering a promising path for technological applications of Kitaev materials.


\vspace{1cm}


\textbf{Keywords:} Quantum matter, topological quantum computation, mean-field theory, Majorana fermions.


\begin{comment}
The Kitaev honeycomb model realizes a quantum spin liquid where spin degrees of freedom are fractionalized into itinerant Majorana fermions and localized vortices.
The entanglement among these particles holds promise for executing fault-tolerant topological quantum computation.
%The entanglement between these particles may be used to perform fault-tolerant topological quantum computation.  Some materials have been proposed to realize a quantum spin liquid phase, but lack strong evidence to confirm the observation of this phase.

While several materials have been proposed to host a quantum spin liquid phase, concrete evidence confirming its existence remains elusive.

%To help identify this phase we propose an electrical test through the charge profile that provides a signature for the $\mathbb{Z}_2$-vortices in Kitaev materials.

To aid in identifying this phase, we propose an electrical test using the charge profile, offering a signature for the $\mathbb{Z}_2$-vortices in Kitaev materials.

%This study explores the electric charge response in Kitaev materials, considering the generalized spin model incorporating Heisenberg and off-diagonal exchange interactions, along with a magnetic Zeeman field of arbitrary directions.
This study explores the electric charge response in Kitaev materials, considering a generalized spin model including Heisenberg, off-diagonal exchange interactions, and magnetic Zeeman field.

Majorana mean-field theory is employed to analyze spin correlations around vortices, revealing spatial anisotropy in the charge distribution, quantified by the electric quadrupole moment.




% This mean-field theory is exact in the pure Kitaev limit and captures the essential physics throughout the Kitaev spin liquid phase.
This mean-field approach remains exact in the pure Kitaev limit and captures the fundamental physics across the Kitaev spin liquid phase.
% The quadrupole-quadrupole interaction between distant vortices is shown to be either repulsive or attractive, depending on the parameters.  We predict that electrically biased scanning probe tips enable the creation of vortices at preselected positions.
The quadrupole-quadrupole interaction between distant vortices is shown to be either repulsive or attractive, depending on the parameters. We predict that electrically biased scanning probe tips can help create vortices at predetermined locations.

The findings of our study pave the way for electrically manipulating Ising anyons in Kitaev spin liquids, offering a promising path for technological applications of Kitaev materials.

% Such findings pave the way for electrically manipulating Ising anyons within Kitaev spin liquids.
% The findings of our study pave the way for electrically manipulating Ising anyons in Kitaev spin liquids, offering a promising path for technological applications of Kitaev materials.

%The techniques studied in this project can also be explored to manipulate the non-Abelian Ising Anyons, providing a path to technological applications for Kitaev materials.


\end{comment}
%Kitaev spin liquids are characterized by the fractionalization of spin-$1/2$ local moments into Majorana fermions and a $\mathbb Z_2$ gauge field, with their properties demonstrated in the exactly solvable Kitaev honeycomb model. Identifying signatures of fractional excitations in Kitaev materials has been a major goal of condensed matter physics, particularly evidenced by phenomena such as the half-quantized thermal Hall conductance observed in $\alpha$-RuCl$_3$. Experimental ambiguity necessitates alternative probes to distinguish Kitaev spin liquids from conventional phases. This study explores the electric charge response in Kitaev materials, considering the generalized spin model incorporating Heisenberg and off-diagonal exchange interactions, along with Zeeman coupling. Majorana mean-field theory is employed to analyze spin correlations around vortices, revealing spatial anisotropy in the charge distribution, quantified by the electric quadrupole moment. potential manipulation of vortices via local electric potentials, offering a promising avenue for electrical detection and manipulation in Kitaev materials, essential for advancing topological quantum computation.
\end{absolutelynopagebreak}


\clearpage

\chapter{Resumo em portugu\^{e}s}
O modelo de Kitaev realiza um l\'iquido qu\^antico de spins onde os graus de liberdade de spin s\~ao fracionados em f\'ermions Majorana itinerantes e v\'ortices localizados. O emaranhamento entre essas part\'iculas é uma rota promissora para a computa\c{c}\~ao qu\^antica topol\'ogica tolerante a falhas. Embora diversos materiais tenham sido propostos para realizar uma fase de l\'iquido qu\^antico de spins, ainda faltam evid\^encias concretas confirmando sua exist\^encias. Para auxiliar na identifica\c{c}\~ao dessa fase, propomos um teste el\'etrico usando o perfil de carga, oferecendo uma assinatura para os v\'ortices $\mathbb{Z}_2$ em materiais de Kitaev.

Este estudo explora a resposta de carga el\'etrica em materiais de Kitaev, considerando o modelo estendido que inclui intera\c{c}\~oes de Heisenberg, $\Gamma$ e campo Zeeman magn\'etico. A teoria de campo m\'edio de Majorana \'e empregada para analisar correla\c{c}\~oes de spin ao redor dos v\'ortices, revelando uma anisotropia espacial na distribui\c{c}\~ao de carga, quantificada pelo momento quadrupolo el\'etrico. Esta abordagem de campo m\'edio permanece exata no limite de Kitaev puro e captura a f\'isica fundamental ao longo da fase  de  l\'iquido de spin de Kitaev. Encontramos  intera\c{c}\~ao quadrupolo-quadrupolo repulsiva ou atrativa  entre v\'ortices distantes, dependendo dos par\^ametros. Prevemos que pontas condutoras polarizadas eletricamente podem criar v\'ortices em locais predefinidos.

Os resultados de nosso estudo abrem caminho para a manipula\c{c}\~ao el\'etrica de anyons de Ising em l\'iquidos de spin de Kitaev, oferecendo um caminho promissor para aplica\c{c}\~oes tecnol\'ogicas de materiais de Kitaev.


\vspace{1cm}


\textbf{Palavras-chave:} Mat\'{e}ria qu\^{a}ntica, computa\c{c}\~{a}o qu\^{a}ntica topol\'{o}gica, teoria de campo m\'{e}dio, f\'{e}rmions de Majorana.


\clearpage


\begin{absolutelynopagebreak}
\chapter{Acknowledgments}


~\vfill


I acknowledge funding by the Brazilian ministry MEC and the Coordenação de Aperfeiçoamento de Pessoal de Nível Superior (CAPES), in special to the Programa Institucional de Internacionalização    PrInt-UFRN which sponsored my one-year exchange study in D\"{u}sseldorf.


\vspace{15mm}
\end{absolutelynopagebreak}






